% =============================================================================
%  Report (EN) — Secondary-neutron dose in water phantoms with monoenergetic photons
%  LaTeX source (article). Compile with pdfLaTeX + BibTeX.
%  Figures are loaded from the figures/ folder via the \smartfig command.
% =============================================================================
\documentclass[11pt,a4paper]{article}

\usepackage[utf8]{inputenc}
\usepackage[T1]{fontenc}
\usepackage[english]{babel}
\usepackage{lmodern}
\usepackage{microtype}
\usepackage[margin=2.3cm]{geometry}

\usepackage{amsmath,amssymb}
\usepackage{graphicx}
\usepackage{booktabs}
\usepackage{caption}
\usepackage{subcaption}
\usepackage{siunitx}
\usepackage{multirow}
\usepackage{authblk}
\usepackage[hidelinks]{hyperref}
\usepackage{cite}

\sisetup{detect-weight=true, detect-family=true}

% --- "Smart" figure: use the PNG if present, otherwise show a placeholder box ---
\newcommand{\smartfig}[3][\linewidth]{%
  \IfFileExists{figures/#2}%
    {\includegraphics[width=#1]{figures/#2}}%
    {\fbox{\parbox[c][0.22\textheight][c]{0.92\linewidth}{\centering\ttfamily
       [missing figure] \detokenize{figures/#2} \\[4pt]\normalfont\small #3}}}%
}

\title{\textbf{Secondary-Neutron Dose in Water Phantoms Irradiated with
Monoenergetic Photons via Monte Carlo Simulation}}

\author[1]{Jorge Ram\'irez L\'opez}
\author[2]{Valeria Catalina Siordia Ort\'iz}
\author[3]{Dr.\ Teodoro Rivera Montalvo (advisor)}
\author[3]{Alejandra Inzunza Lugo (advisor)}
\author[3]{C\'esar Samuel Romero N\'u\~nez (advisor)}
\affil[1]{Physics B.Sc., Division of Basic Sciences, CUCEI, University of Guadalajara \\ \texttt{jorge.ramirez9331@alumnos.udg.mx}}
\affil[2]{Biomedical Engineering, Tlajomulco University Center, University of Guadalajara}
\affil[3]{Dosimetry Laboratory, CICATA Legaria Unit, Instituto Polit\'ecnico Nacional}
\date{}

\begin{document}
\maketitle

\begin{abstract}
\noindent
Understanding the spatial dose distribution is essential to assess how beam energy
reaches the tumor volume while minimizing exposure of surrounding healthy tissue. In
particular, secondary neutrons produced by photonuclear processes are an additional
radiation source that, although small in magnitude, can affect out-of-field dose and
raise the risk of unwanted biological effects. In this work we characterized the
spatial dose distribution in a water phantom using Monte Carlo simulations in
\texttt{Geant4} for monoenergetic photon beams of 6, 10, 12, 15 and 18~MeV. For each
energy, 100 runs of \(10^{6}\) particles were performed, and the results were
processed in \texttt{MATLAB} to obtain percentage depth-dose (PDD) curves, global
metrics (\(D_{\max}\) and its depth), 2D/3D maps per component and energy histograms.
The total dose was decomposed into primary and secondary contributions (electrons,
positrons, photons and \textbf{neutrons}), emphasizing the neutron component. The PDD
curves reproduced the expected physics. A low-magnitude neutron component was
identified, growing with energy and spatially more diffuse, consistent with
photonuclear production. This highlights the importance of quantifying the neutron
component for its potential impact on cancer induction and the need to account for it
in dosimetric studies and radiotherapy treatment planning.
\end{abstract}

\noindent\textbf{Keywords:} Monte Carlo, Geant4, PDD, photoneutrons, dosimetry.

% =============================================================================
\section{Introduction}
Cancer, the abnormal growth of malignant cells, can be treated with radiotherapy: the
use of ionizing radiation to treat localized tumors. Its application demands accurate
dose characterization, both inside and outside the treatment
field~\cite{khan2014radiation}. The percentage depth dose (PDD) describes how
absorbed dose varies with depth when a beam is incident on a patient or
phantom~\cite{khan2014radiation}. To this end we used Monte Carlo simulations---the
gold standard~\cite{fielding2023montecarlo}---with the Geant4
toolkit~\cite{agostinelli2003geant4}, which accurately models radiation--matter
interactions with full traceability.

In high-energy photon treatments (\(\gtrsim 10~\mathrm{MV}\)), photonuclear reactions
can occur and generate a field of secondary neutrons (thermal and fast) that
contaminates the therapeutic beam with unwanted dose~\cite{banaee2021}. Its impact is
significant out of field, where commercial treatment planning systems tend to
underestimate this component~\cite{kry2017tg158}. Quantifying the neutron
contribution matters because of its Relative Biological Effectiveness
(RBE)~\cite{johns1983physics}: neutrons peak in effectiveness for complex DNA damage
near \(1~\mathrm{MeV}\)~\cite{baiocco2016origin}, so even at low doses their toxicity
motivates characterizing the physical dose they deposit~\cite{kry2017tg158}.

Despite this evidence, Mexican regulation shows a gap: NOM-033-NUCL-2016
\cite{nom033_2016} requires neutron detectors but sets no operational procedures for
significant measurements, and neither it nor NOM-002-SSA3-2017~\cite{nom002ssa32017}
details how to integrate neutron dose into clinical planning. This work evaluates the
spatial dose distribution of monoenergetic photon beams (6, 10, 12, 15 and 18~MeV) in
a voxelized water phantom, decomposing the dose into primary and secondary components
with emphasis on the neutron component. Processing was done in MATLAB R2025a on
high-performance computing infrastructure, within the XXX Delf\'in Program research
summer~\cite{programadelfin2025} at CICATA-Legaria, IPN.

% =============================================================================
\section{Theoretical background}
\textbf{Absorbed dose} \(D\) is the energy deposited per unit mass
(\(D=\mathrm{d}\epsilon/\mathrm{d}m\)); its SI unit is the gray
(Gy)~\cite{attix2004dosimetry}. The \textbf{percentage depth dose} normalizes the dose
at depth \(d\) to a reference depth (typically \(d_{\max}\))~\cite{khan2014radiation}:
\begin{equation}
  \mathrm{PDD}(d,\mathrm{SSD}) = 100 \times
  \frac{D(d,\mathrm{SSD})}{D(d_{\mathrm{ref}},\mathrm{SSD})}.
\end{equation}

Above \(\sim 10~\mathrm{MeV}\), photons can produce neutrons via \textbf{photonuclear
reactions}, a process dominated by the Giant Dipole Resonance; for oxygen-16 in water
the cross section peaks near 23~MeV~\cite{geant4physref}. Photoneutrons thermalize
through collisions with hydrogen down to \(\sim
0.025~\mathrm{eV}\)~\cite{lakshmanan1982}. Since dose alone does not predict
biological effect, the RBE compares a test radiation to a reference for the same
effect~\cite{reeve2018}. Neutrons produce high linear-energy-transfer (LET)
secondaries, causing complex damage and an RBE significantly greater than
1~\cite{reeve2018}.

% =============================================================================
\section{Methodology}

\subsection{Geant4 configuration}
Simulations ran in \texttt{Geant4 v11.3.2} in multithreaded mode. Electromagnetic
physics used \texttt{G4EmStandardPhysics}; neutron production by photonuclear
reactions was enabled with \texttt{G4EmExtraPhysics} (\texttt{GammaNuclear} process)
together with \texttt{G4HadronPhysicsQGSP\_BIC\_HP} and
\texttt{G4HadronElasticPhysics} for elastic scattering and inelastic neutron
interactions with thermal-to-intermediate precision (High Precision model). This is
an optimized version of the \texttt{QGSP\_BIC\_HP} list~\cite{geant4qgspbic}, without
decay or ion processes given their low relevance for 6--18~MeV photon beams.

\subsection{Simulation environment}
The environment was implemented with \texttt{G4VUserDetectorConstruction}, defining a
cubic water phantom (\texttt{G4\_WATER}) 40~cm on a side. The volume was discretized
into a \(27\times27\times200\) voxel grid; since 40~cm is not divisible by the nominal
15~mm, the real pitch is \(400/27=14.8148\)~mm in \(X/Y\) and exactly 2~mm in \(Z\)
(volume \(\approx 439~\mathrm{mm}^3\) per voxel).\footnote{The nominal design was
\(15\times15\times2~\mathrm{mm}^3\); rounding to 27 voxels fixes the quoted real
pitch.} A geometry sketch is shown in Fig.~\ref{fig:esquema} and a view of the
environment in Fig.~\ref{fig:entorno}.

\begin{figure}[t]
  \centering
  \smartfig[0.8\linewidth]{esquema_geometrico.png}{Geometry sketch.}
  \caption{Geometry sketch of the Geant4 simulation environment.}
  \label{fig:esquema}
\end{figure}

\begin{figure}[t]
  \centering
  \smartfig[0.85\linewidth]{visualizador_18MeV.png}{Beam tracks in the phantom.}
  \caption{Tracks of a narrow beam of 10 primary 18~MeV photons interacting with the
  water phantom. The mesh shown is low-resolution and illustrative only.}
  \label{fig:entorno}
\end{figure}

\subsection{Beam generation and execution}
The beam was configured with \texttt{G4VUserPrimaryGeneratorAction} and controlled via
macros: monoenergetic photons of 6, 10, 12, 15 and 18~MeV incident perpendicularly
from a centered point source (ideal flat field). Per-run actions used
\texttt{G4UserRunAction} and \texttt{G4VUserActionInitialization}, recording the
energy deposited by the primary beam and by secondaries. Data were saved to
\texttt{.csv} (one per run) with all non-zero-dose voxels and the per-component
breakdown. Each simulation used \(10^{6}\) particles per run, with 100 independent
runs per energy. Batch execution was automated with \texttt{bash} and macro files.

\subsection{Processing and uncertainty estimation}
Processing was automated in \texttt{MATLAB}~\cite{matlab2024}: reading the raw CSVs,
computing the 3D mean dose per voxel, its standard deviation and the PDD curve;
per-run and per-energy metrics (\(D_{\max}\), its depth, means, standard deviations
and \(\epsilon=100\,\sigma/\mu\)); histograms; 2D/3D maps per component; and export of
averaged files and summary tables.

For each voxel with mean dose \(\bar D>0\), the statistical uncertainty of the mean
was estimated as \(\epsilon_{\text{vox}} =
100\,(\sigma_{\text{vox}}/\sqrt{N})/\bar D\), reporting its spatial average
\(\overline{\epsilon}_{\text{vox}}\) over voxels with \(\bar D>0\). Using the relative
SEM is consistent with practice in established Monte Carlo
codes~\cite{mcnp_v4c2000}. For global per-energy metrics we report
\(\epsilon=100\,\sigma/\mu\), quantifying run-to-run variability. Both characterize
complementary aspects of the uncertainty.

% =============================================================================
\section{Results}

\begin{figure}[t]
  \centering
  \smartfig[0.85\linewidth]{PDD_comparacion.png}{PDD curves per energy.}
  \caption{PDD curves for the considered energies.}
  \label{fig:pdd}
\end{figure}

\begin{table}[t]
  \centering
  \caption{Summary of dosimetric parameters per beam energy. Means \(\mu\) and
  relative errors \(\epsilon=100\,\sigma/\mu\) estimated from 100 runs per energy.}
  \label{tab:resumen}
  \begin{tabular}{l S[table-format=2.2] S[table-format=1.2] S[table-format=2.2] S[table-format=2.2] S[table-format=1.2] S[table-format=1.2]}
    \toprule
    \multirow{2}{*}{\textbf{Energy}} &
      \multicolumn{2}{c}{\textbf{Max.\ dose}} &
      \multicolumn{2}{c}{\textbf{Depth of \(D_{\max}\)}} &
      \multicolumn{2}{c}{\textbf{Total dose}} \\
    \cmidrule(lr){2-3}\cmidrule(lr){4-5}\cmidrule(l){6-7}
      & {\(\mu\) [\si{\nano\gray}]} & {\(\epsilon\) [\%]}
      & {\(\mu\) [\si{\milli\metre}]} & {\(\epsilon\) [\%]}
      & {\(\mu\) [\si{\micro\gray}]} & {\(\epsilon\) [\%]} \\
    \midrule
    6 MeV  & 11.03 & 0.59 & 26.26 & 10.10 & 1.61 & 0.08 \\
    10 MeV & 15.59 & 0.50 & 43.88 &  6.47 & 2.37 & 0.09 \\
    12 MeV & 17.79 & 0.42 & 52.58 &  6.12 & 2.74 & 0.11 \\
    15 MeV & 21.12 & 0.35 & 65.72 &  4.17 & 3.28 & 0.12 \\
    18 MeV & 24.39 & 0.39 & 77.26 &  3.94 & 3.82 & 0.11 \\
    \bottomrule
  \end{tabular}
\end{table}

\begin{table}[t]
  \centering
  \caption{Mean per-voxel relative error of the mean dose
  (\(\overline{\epsilon}_{\text{vox}}=\mathrm{SEM}/\bar D\), in \%). Spatial average
  over voxels with dose \(>0\).}
  \label{tab:err_voxel}
  \begin{tabular}{l S[table-format=1.2]}
    \toprule
    \textbf{Energy} & {\(\overline{\epsilon}_{\text{vox}}\) [\%]} \\
    \midrule
    6 MeV  & 5.53 \\
    10 MeV & 5.54 \\
    12 MeV & 5.49 \\
    15 MeV & 5.44 \\
    18 MeV & 5.36 \\
    \bottomrule
  \end{tabular}
\end{table}

\begin{figure}[t]
  \centering
  \begin{subfigure}[t]{0.32\textwidth}\centering
    \smartfig{Hist_DosisMax_6MeV.png}{6 MeV}\caption{6 MeV}\end{subfigure}\hfill
  \begin{subfigure}[t]{0.32\textwidth}\centering
    \smartfig{Hist_DosisMax_10MeV.png}{10 MeV}\caption{10 MeV}\end{subfigure}\hfill
  \begin{subfigure}[t]{0.32\textwidth}\centering
    \smartfig{Hist_DosisMax_12MeV.png}{12 MeV}\caption{12 MeV}\end{subfigure}
  \\[4pt]
  \begin{subfigure}[t]{0.32\textwidth}\centering
    \smartfig{Hist_DosisMax_15MeV.png}{15 MeV}\caption{15 MeV}\end{subfigure}\hfill
  \begin{subfigure}[t]{0.32\textwidth}\centering
    \smartfig{Hist_DosisMax_18MeV.png}{18 MeV}\caption{18 MeV}\end{subfigure}\hfill
  \begin{subfigure}[t]{0.32\textwidth}\end{subfigure}
  \caption{Histograms of \textbf{maximum dose} (nGy) per energy; \(\mu\), \(\sigma\)
  and \(\epsilon\) values in Table~\ref{tab:resumen}.}
  \label{fig:hist-dosis}
\end{figure}

\begin{figure}[t]
  \centering
  \smartfig[0.85\linewidth]{Maps2D_15MeV.png}{2D maps at 15 MeV.}
  \caption{2D dose maps (maximum-intensity projection in XY) for the 15~MeV beam:
  total dose, primary photons, electrons, positrons, photons, neutrons and other
  secondaries.}
  \label{fig:maps2d}
\end{figure}

\begin{figure}[t]
  \centering
  \begin{subfigure}[b]{0.32\textwidth}\centering
    \smartfig{Map3D_15MeV_Total.png}{Total}\caption{Total dose}\end{subfigure}\hfill
  \begin{subfigure}[b]{0.32\textwidth}\centering
    \smartfig{Map3D_15MeV_eminus.png}{e-}\caption{Electrons}\end{subfigure}\hfill
  \begin{subfigure}[b]{0.32\textwidth}\centering
    \smartfig{Map3D_15MeV_eplus.png}{e+}\caption{Positrons}\end{subfigure}
  \\[4pt]
  \begin{subfigure}[b]{0.32\textwidth}\centering
    \smartfig{Map3D_15MeV_neutron.png}{n}\caption{Neutrons}\end{subfigure}\hfill
  \begin{subfigure}[b]{0.32\textwidth}\centering
    \smartfig{Map3D_15MeV_other.png}{other}\caption{Other}\end{subfigure}\hfill
  \begin{subfigure}[b]{0.32\textwidth}\end{subfigure}
  \caption{3D dose distribution for the 15~MeV beam. Voxels are rendered if their
  normalized dose exceeds 25\,\% (1\,\% for neutrons).}
  \label{fig:maps3d}
\end{figure}

% =============================================================================
\section{Conclusions}
The PDD curves (Fig.~\ref{fig:pdd}) reproduce the expected physics: build-up region,
\(D_{\max}\) depth increasing with energy, and a soft tail, supporting model
consistency. Global metrics are highly stable: \(D_{\max}\) shows very low relative
errors (\(\epsilon\approx 0.35\text{--}0.59\%\)), while the depth of \(D_{\max}\) has
larger dispersion that decreases with energy (\(\epsilon\) from \(\sim 10\%\) to
\(\sim 4\%\)), in agreement with the histograms (Fig.~\ref{fig:hist-dosis}) and
Table~\ref{tab:resumen}. Spatially, the per-voxel relative SEM stayed within
\(5.3\text{--}5.5\%\) (Table~\ref{tab:err_voxel}).

The 2D/3D visualizations (Figs.~\ref{fig:maps2d} and \ref{fig:maps3d}) show that
charged components (\(e^-/e^+\)) concentrate dose near the beam axis, whereas
secondary neutrons appear more scattered and at low levels, with greater relative
presence at high energies. Although their dosimetric contribution is small, their
potential radiobiological relevance justifies characterizing them separately and,
in future work, translating them via LET and RBE. This work provides a reproducible
framework to compare energies and dose components.

\subsection*{Limitations and future work}
(1) Monoenergetic, idealized beams; (2) homogeneous phantom without accelerator head;
(3) fixed voxel resolution; (4) LET/RBE not evaluated. Proposed extensions:
incorporate clinical spectra and head geometry; study out-of-field scenarios; compute
LET and estimate RBE for secondaries (especially neutrons); and apply
variance reduction to lower \(\overline{\epsilon}_{\text{vox}}\) at fixed compute cost.

\subsection*{Acknowledgments}
To the Delf\'in Program and fellow participants; to CICATA-Legaria and the Dosimetry
Laboratory team; and especially to Dr.\ Teodoro, Alejandra and C\'esar for their
guidance, patience and support.

\bibliographystyle{unsrt}
\bibliography{references}

\end{document}
